% Part III-A, results of the offline code judge study (included by main_en.tex).

\section{What the judges see in code}\label{sec:res-study}

\paragraph{The visible tests leave a large gap.} The base policy's \NUM{cs_policy} single attempts (\NUM{cs_samples}
per problem) pass the visible tests in \NUM{cs_vispass_all} of cases and every test in \NUM{cs_allpass_all}. Of the
attempts that pass the visible tests, \NUM{cs_hack_all} fail a held-out test: \NUM{cs_hack_kod} on KodCode val,
\NUM{cs_hack_mbpp} on MBPP+ and \NUM{cs_hack_he} on HumanEval+. A reward built from the visible tests pays these
attempts as much as the correct ones, before any training has had a chance to exploit it.

\paragraph{Both judges see most of it, and \kev{} is no better on the training distribution.}
Table~\ref{tab:code-study} gives, among the attempts that pass the visible tests, each signal's AUROC for also passing
the held-out tests. Over all three sets, \kev{}'s \emph{genuine} reaches \NUM{csauc_all_kevg}
\ci{\NUM{csauclo_all_kevg}}{\NUM{csauchi_all_kevg}} and the AR judge's \NUM{csauc_all_arg}
\ci{\NUM{csauclo_all_arg}}{\NUM{csauchi_all_arg}} (95\% intervals from resampling problems). The difference,
\NUM{csd_all_kevgarg} \ci{\NUM{csdlo_all_kevgarg}}{\NUM{csdhi_all_kevgarg}}, is not resolved. With the \emph{correct}
question \kev{} is ahead by \NUM{csd_all_kevcarc} \ci{\NUM{csdlo_all_kevcarc}}{\NUM{csdhi_all_kevcarc}}. Most of
either difference comes from MBPP+, where \kev{} leads by \NUM{csd_mbpp_kevgarg}
\ci{\NUM{csdlo_mbpp_kevgarg}}{\NUM{csdhi_mbpp_kevgarg}} on \emph{genuine}. On KodCode val, the distribution the policy
is trained on, the AR judge is ahead by a margin that is not resolved (\NUM{csd_kod_kevgarg}
\ci{\NUM{csdlo_kod_kevgarg}}{\NUM{csdhi_kod_kevgarg}}), and on HumanEval+ the two cannot be told apart
(\NUM{csd_he_kevgarg} \ci{\NUM{csdlo_he_kevgarg}}{\NUM{csdhi_he_kevgarg}}). Over all attempts, both judges separate the ones that pass
every test from the rest better (\NUM{csaucall_all_kevg} and \NUM{csaucall_all_arg}), with the AR judge slightly ahead
(\NUM{csdall_all_kevgarg} \ci{\NUM{csdalllo_all_kevgarg}}{\NUM{csdallhi_all_kevgarg}}).

\begin{table}[t]
\centering\small
\begin{tabularx}{\linewidth}{@{}>{\raggedright\arraybackslash}Xcccccc@{}}
\toprule
 & \multicolumn{5}{c}{held-out tests pass, among visible-pass attempts} & all pass \\ \cmidrule(lr){2-6}\cmidrule(l){7-7}
 & KodCode val & MBPP+ & HumanEval+ & all [95\% CI] & within & all \\ \midrule
attempts passing the visible tests & 391 & 225 & 116 & 732 & & \\
\quad of which fail a held-out test & 169 & 57 & 37 & 263 & & \\ \midrule
\kev{} \emph{genuine} & 0.89 & 0.83 & 0.85 & 0.87 \ci{0.84}{0.90} & 0.81 & 0.93 \\
\kev{} \emph{correct} & 0.88 & 0.82 & 0.86 & 0.87 \ci{0.84}{0.90} & 0.79 & 0.93 \\
AR \emph{genuine} & 0.90 & 0.74 & 0.84 & 0.86 \ci{0.83}{0.89} & 0.85 & 0.94 \\
AR \emph{correct} & 0.87 & 0.74 & 0.84 & 0.84 \ci{0.81}{0.87} & 0.81 & 0.93 \\
\bottomrule
\end{tabularx}

\caption{AUROC of each judge signal for passing the held-out tests, among the base policy's attempts that pass the
visible tests (\NUM{cs_samples} attempts per problem at temperature 1). \emph{within}: mean AUROC within the
\NUM{cs_mixed_all} problems that have visible-pass attempts both passing and failing the held-out tests. \emph{all
pass}: AUROC for passing every test, over all attempts. 95\% intervals from resampling problems.}
\label{tab:code-study}
\end{table}

\paragraph{Within a problem the evidence is thin.} GRPO's group baseline compares the episodes of one problem, so the
judge factor in \texttt{vis\_kev} matters through how it ranks a problem's visible-pass episodes. Only
\NUM{cs_mixed_all} of the \NUM{cs_problems_all} problems with a visible-pass attempt have attempts on both sides of the
held-out tests. Within them \kev{}'s \emph{genuine} reaches \NUM{csaucw_all_kevg} and the AR judge's
\NUM{csaucw_all_arg}, a difference of \NUM{csdw_all_kevgarg}
\ci{\NUM{csdwlo_all_kevgarg}}{\NUM{csdwhi_all_kevgarg}}. These problems are too few to rank the judges, and too few to
predict how much the factor will change training.

\paragraph{Hard-coded lookups are caught.} \NUM{cs_hardcoded} lookups of the visible tests' calls pass the visible tests
of their MBPP+ or HumanEval+ problem (\NUM{cs_hardcoded_heldout} of them also passes the held-out tests). Both judges
separate them from the reference solutions of the same problems almost perfectly (Table~\ref{tab:code-hardcoded}).
\kev{}'s \emph{genuine} gives a lookup \NUM{rh_hc_kevg} on average and a reference solution \NUM{rh_ref_kevg}, and never
scores a lookup above the reference solution of its own problem. Such lookups are the crudest way to special-case the
tests. That the judges catch them shows that the \emph{genuine} question is understood. It does not show that subtler
over-fitting would be caught.

\begin{table}[t]
\centering\small
\begin{tabular}{@{}lcccc@{}}
\toprule
 & & \multicolumn{2}{c}{mean $p(\text{yes})$} & lookup scored \\ \cmidrule(lr){3-4}
signal & AUROC & reference & lookup & above reference \\ \midrule
\kev{} \emph{genuine} & 1.000 & 0.75 & 0.03 & 0.0\% \\
\kev{} \emph{correct} & 0.997 & 0.75 & 0.07 & 0.2\% \\
AR \emph{genuine} & 1.000 & 0.80 & 0.02 & 0.0\% \\
AR \emph{correct} & 0.991 & 0.79 & 0.05 & 0.8\% \\
\bottomrule
\end{tabular}

\caption{Reference solutions against hard-coded lookups of the visible tests, on the \NUM{rh_problems} MBPP+ and
HumanEval+ problems whose visible tests allow a lookup: AUROC, mean $p(\text{yes})$ of each kind, and how often the
lookup scores above its own problem's reference solution.}
\label{tab:code-hardcoded}
\end{table}

\paragraph{What this predicts.} A judge factor can down-weight most of the visible-pass attempts that are wrong, and
\kev{} does that about as well as the AR judge, a general model of the same size asked the same question. Whether the factor
reduces that gap during training, and at what cost to the visible-test pass rate, is what the training runs measure.
